%%% define the colour of the title				%%%
%%% could be set to match general colour theme 	%%%

\tikzset{external/export=false}
\begin{titlepage}
    %\color{bordeau} \hfill \vfill \tiny \ecodocnum

    \begin{textblock}{5}(0,0)
      \textblockcolour{bordeau}
      %\vspace{10mm}
      \includegraphics [scale=1]{00_Guards/media/bande.png}
      \vspace{300mm}
    \end{textblock}
    
    \begin{textblock}{1}(0.6,9.5)
      \Huge{\rotatebox{90}{\color{white}{\fontsize{38}{54}\selectfont Thèse de doctorat}}}
    \end{textblock}
    
    \begin{textblock}{1}(0.6,3)
      \Large{\rotatebox{90}{\color{white}{NNT : \NNT}}}
    \end{textblock}
    
    \begin{textblock}{1}(\hpostt,\vpostt)
      \textblockcolour{white}
      \includegraphics[scale=1]{00_Guards/media/etab/\logoEtt.png}
    \end{textblock}
    
    \begin{textblock}{1}(\hpos,\vpos)
      \textblockcolour{white}
      \includegraphics[scale=1]{00_Guards/media/etab/\logoEt.png}
    \end{textblock}
    
    %\vspace{6cm}
    %% Texte
    \begin{textblock}{10}(5.7,3)
      \textblockcolour{white}
      
      \color{bordeau}
      %\begin{center}  
      \begin{flushright}
        \huge{\PhDTitle} \bigskip %% Titre de la thèse 
        \vfill
        \color{black} %% Couleur noire du reste du texte
        \normalsize {Thèse de doctorat de l'Université Paris-Saclay} \\
        préparée à \PhDworkingplace \\ \bigskip
        \vfill
        École doctorale n$^{\circ}$\ecodocnum ~\ecodoctitle ~(\ecodocacro)  \\
        
        \small{Spécialité de doctorat: \PhDspeciality} \bigskip %% Spécialité 
        \vfill  
        \footnotesize{Thèse présentée et soutenue à \defenseplace, le \defensedate, par} \bigskip
        \vfill
        \Large{\textbf{\textsc{\PhDname}}} %% Nom du docteur
        \vfill
        %\bigskip
      \end{flushright}
      
      %\end{center}
      \color{black}
      %% Jury
      \begin{flushleft}

        \small Composition du Jury :
      \end{flushleft}
      %% Members of the jury
      
      \small
      %\begin{center}
      \newcolumntype{L}[1]{>{\raggedright\let\newline\\\arraybackslash\hspace{0pt}}m{#1}}
      \newcolumntype{R}[1]{>{\raggedleft\let\newline\\\arraybackslash\hspace{0pt}}lm{#1}}
      
      \label{jury} 																				%% Mettre à jour si des membres ont été ajoutés ou retirés / Update if members have been added or removed
      \begin{flushleft}
        \begin{tabular}{@{} L{9.5cm} R{4.5cm}}
          \jurynameA \\ \juryadressA & \juryroleA \\[5pt]
          \jurynameB \\ \juryadressB & \juryroleB \\[5pt]
          \jurynameC \\ \juryadressC & \juryroleC \\[5pt]
          \jurynameD \\ \juryadressD & \juryroleD \\[5pt]
          \jurynameE \\ \juryadressE & \juryroleE \\[5pt]
          \jurynameF \\ \juryadressF & \juryroleF \\[5pt]
          \jurynameG \\ \juryadressG & \juryroleG \\[5pt]
          \jurynameH \\ \juryadressH & \juryroleH \\[5pt]
        \end{tabular}
      \end{flushleft}
      %\end{center}
    \end{textblock}
\end{titlepage}